\documentclass[10pt]{article}
\providecommand{\FIGDIR}{fig}
\newcommand{\DIGESTDATE}{11 September 2026}
% ===================== PM Research Watch - shared preamble =====================
\usepackage[T1]{fontenc}
\usepackage[utf8]{inputenc}
\usepackage{amsmath,amssymb}
\usepackage{tgpagella}
\usepackage{tgheros}
\renewcommand{\sfdefault}{qhv}
\renewcommand{\ttdefault}{lmtt}
\usepackage{microtype}
\usepackage[english]{babel}

\usepackage[a4paper,top=17mm,bottom=17mm,left=17mm,right=17mm,
            headsep=6mm,footskip=10mm]{geometry}
\usepackage{graphicx}
\usepackage{xcolor}
\usepackage{colortbl}
\usepackage{tikz}
\usetikzlibrary{calc}
\usepackage[most]{tcolorbox}
\usepackage{titlesec}
\usepackage{fancyhdr}
\usepackage{booktabs}
\usepackage{array}
\usepackage{longtable}
\usepackage{enumitem}
\usepackage{multicol}
\usepackage{ragged2e}
\usepackage[hidelinks]{hyperref}

% ------------------------------------------------------------------ palette
\definecolor{Ink}    {HTML}{14262E}
\definecolor{Deep}   {HTML}{0F4C5C}
\definecolor{Teal}   {HTML}{1F7A8C}
\definecolor{Sky}    {HTML}{4FA3C4}
\definecolor{Sage}   {HTML}{6FA287}
\definecolor{Amber}  {HTML}{E3A008}
\definecolor{Clay}   {HTML}{B5643C}
\definecolor{Coral}  {HTML}{D1495B}
\definecolor{Violet} {HTML}{6A5B9E}
\definecolor{Slate}  {HTML}{4A6572}
\definecolor{Paper}  {HTML}{FBF9F5}
\definecolor{Mist}   {HTML}{EEEAE2}
\definecolor{Rule}   {HTML}{D9D3C9}

\pagecolor{Paper}
\color{Ink}

\hypersetup{colorlinks=true, linkcolor=Deep, urlcolor=Teal, citecolor=Teal,
            pdftitle={PM Research Watch}, pdfauthor={Automated daily digest}}

% ------------------------------------------------------------------ headings
\titleformat{\section}
  {\sffamily\bfseries\large\color{Deep}}{}{0pt}
  {}[{\vspace{1.2mm}\color{Amber}\titlerule[1.1pt]}]
\titlespacing*{\section}{0pt}{5.5mm}{2.8mm}

\titleformat{\subsection}
  {\sffamily\bfseries\normalsize\color{Teal}}{}{0pt}{}
\titlespacing*{\subsection}{0pt}{3.6mm}{1.4mm}

\setlist[itemize]{leftmargin=4.2mm, itemsep=1.1mm, topsep=1.3mm, parsep=0pt,
                  label=\textcolor{Teal}{\footnotesize$\blacktriangleright$}}

% ------------------------------------------------------------------ header/footer
\pagestyle{fancy}
\fancyhf{}
\renewcommand{\headrulewidth}{0pt}
\renewcommand{\footrulewidth}{0pt}
\fancyhead[L]{\sffamily\scriptsize\color{Slate}PM RESEARCH WATCH\;\textcolor{Rule}{\textbar}\;\DIGESTDATE}
\fancyhead[R]{\sffamily\scriptsize\color{Slate}Particulate matter monitoring \& health effects}
\fancyfoot[C]{\sffamily\scriptsize\color{Slate}\thepage}

% ------------------------------------------------------------------ components
% metric tile
\newcommand{\metric}[3]{%
  \begin{tcolorbox}[enhanced, colback=white, colframe=Rule, boxrule=0.5pt,
    arc=1.6mm, left=2.2mm, right=2.2mm, top=1.8mm, bottom=1.8mm,
    borderline west={1.7mm}{0pt}{#3}, nobeforeafter, valign=center,
    height=17mm, width=\linewidth]
    \RaggedRight
    {\sffamily\bfseries\fontsize{17}{17}\selectfont\color{#3}#1}\\[0.5mm]
    {\sffamily\fontsize{6.5}{7.4}\selectfont\color{Slate}#2}
  \end{tcolorbox}}

% pull-quote / key insight box
\newtcolorbox{keybox}[1][]{enhanced, colback=Mist, colframe=Mist,
  arc=1.6mm, left=3.2mm, right=3.2mm, top=2.4mm, bottom=2.4mm,
  borderline west={1.7mm}{0pt}{Amber}, #1}

% subtopic band
\newcommand{\band}[2]{%
  \vspace{3.2mm}\noindent
  \begin{tcolorbox}[enhanced, colback=#2, colframe=#2, arc=1.2mm,
    left=2.6mm, right=2.6mm, top=1.1mm, bottom=1.1mm, boxrule=0pt,
    width=\linewidth, nobeforeafter]
    {\sffamily\bfseries\small\color{white}#1}
  \end{tcolorbox}\vspace{1.6mm}}

% one paper entry: {short cite}{title}{journal}{doi}{body}
\newcommand{\paperentry}[5]{%
  \par\vspace{1.9mm}\noindent
  {\sffamily\bfseries\footnotesize\color{Ink}#1}\;%
  {\sffamily\fontsize{7}{8}\selectfont\color{Slate}#3}\\[0.4mm]
  {\footnotesize\itshape\color{Deep}#2}\\[0.6mm]
  {\footnotesize #4}\\[0.3mm]
  {\fontsize{6.6}{7.6}\selectfont\color{Slate}\texttt{doi:}\,\href{https://doi.org/#5}{\texttt{\detokenize{#5}}}}
  \par}

% caption style
\newcommand{\figcap}[1]{\par\vspace{1.0mm}
  {\fontsize{7.2}{8.6}\selectfont\color{Slate}\textsf{\textbf{Fig.}}\; #1\par}\vspace{1.5mm}}

\providecommand{\permil}{\textperthousand}

\setlength{\parindent}{0pt}
\setlength{\parskip}{0pt}
\renewcommand{\arraystretch}{1.18}


\begin{document}

% ==================================================================== MASTHEAD
\thispagestyle{empty}
\begin{tikzpicture}[remember picture, overlay]
  \fill[Deep] (current page.north west) rectangle
      ($(current page.north east)+(0,-67mm)$);
  \fill[Amber] ($(current page.north west)+(0,-67mm)$) rectangle
      ($(current page.north east)+(0,-68.4mm)$);
  % subtle particle field
  \foreach \x/\y/\r in {22/12/0.9, 41/26/1.5, 63/9/0.7, 88/31/1.1, 112/17/1.9,
                        138/29/0.8, 159/13/1.3, 178/54/1.0, 196/21/1.6,
                        30/44/0.6, 74/57/1.2, 125/48/0.9, 168/6/0.8, 191/59/0.7,
                        52/18/1.0, 100/60/0.6, 148/20/1.1, 66/49/1.4, 205/45/0.8}
    \fill[Sky, opacity=0.30] ($(current page.north west)+(\x mm,-\y mm)$) circle (\r mm);
\end{tikzpicture}

\vspace*{4mm}
{\sffamily\fontsize{7.6}{9}\selectfont\color{Sky}%
 AUTOMATED DAILY LITERATURE DIGEST\;\textbullet\;PhD RESEARCH WATCH\par}
\vspace{2.2mm}
{\sffamily\bfseries\fontsize{27}{29}\selectfont\color{white}%
 Particulate Matter\par}
{\sffamily\fontsize{27}{29}\selectfont\color{Sky}%
 Monitoring \& Health Effects\par}
\vspace{5.0mm}
{\sffamily\fontsize{8.4}{10}\selectfont\color{white}%
 Issue 2026-09-11\;\;\textcolor{Amber}{\textbar}\;\;Entry window 10--11 September 2026\;\;%
 \textcolor{Amber}{\textbar}\;\;10 records screened in scope\par}
\vspace{18mm}

% -------------------------------------------------------------- metric strip
\noindent
\begin{minipage}[t]{0.190\linewidth}\metric{10}{RECORDS IN SCOPE\\(28 EXCLUDED)}{Deep}\end{minipage}\hfill
\begin{minipage}[t]{0.190\linewidth}\metric{8}{SUBTOPIC CLUSTERS\\REPRESENTED}{Teal}\end{minipage}\hfill
\begin{minipage}[t]{0.190\linewidth}\metric{4}{EFFECT ESTIMATES\\EXTRACTED}{Sky}\end{minipage}\hfill
\begin{minipage}[t]{0.190\linewidth}\metric{2}{SENSING \&\\INSTRUMENTATION}{Amber}\end{minipage}\hfill
\begin{minipage}[t]{0.190\linewidth}\metric{3}{METADATA ONLY\\(NO ABSTRACT)}{Clay}\end{minipage}

\vspace{6mm}

% -------------------------------------------------------------- signal of day
\section*{\sffamily\bfseries\color{Deep}Signal of the day}
\vspace{-1.5mm}{\color{Amber}\rule{\linewidth}{1.1pt}}\vspace{2.5mm}

\begin{keybox}
{\sffamily\bfseries\small\color{Deep}Eight million people, intervals three decimal places
wide --- and the number that decides what the study means is the one nobody
measured.}\\[1.4mm]
\footnotesize
Lee et al.\ (\textit{Allergy}) linked 8{,}252{,}498 Korean National Health Insurance
enrollees to 1\,km long-term PM estimates and got incident-asthma hazard ratios of
\textbf{1.33 (95\% CI 1.32--1.34)} per 10\,$\mu$g\,m$^{-3}$ of PM$_{2.5}$ and
\textbf{1.38 (1.37--1.38)} for PM$_{10}$. Stratifying on age, Medical Aid status and
urbanicity produced a clean monotone gradient from 1.16 (1.14--1.18) in younger rural
adults without Medical Aid to \textbf{1.80 (1.66--1.94)} in older urban Medical Aid
recipients, interaction $p<0.001$. At this sample size the confidence intervals stop being
the interesting uncertainty. Incident asthma here is a diagnosis code, and the same
gradient that defines the exposure strata --- urban, poor, old --- is the gradient of
healthcare contact, so a portion of the 1.16-to-1.80 spread is ascertainment rather than
susceptibility, and the paper offers no lever to separate them. The second problem is
ours. The modifier being tested is \emph{urbanicity}, a within-city contrast, and the
exposure surface resolving it is 1\,km. Intra-urban PM$_{2.5}$ gradients run at hundreds of
metres near roads; a 1\,km surface averages the Medical Aid recipient beside the arterial
together with the resident four blocks back. That kind of misclassification does not
manufacture a gradient, it compresses one --- so 1.80 is probably a floor, and the study's
own design is an argument for finer exposure surfaces.
\end{keybox}

\vspace{3mm}

% -------------------------------------------------------------- exec brief
\section*{\sffamily\bfseries\color{Deep}What changed today}
\vspace{-1.5mm}{\color{Amber}\rule{\linewidth}{1.1pt}}\vspace{2.5mm}

\footnotesize
\begin{itemize}
\item \textbf{A controlled human exposure asks whether blood can stand in for a
bronchoscope.} Thirty-two participants, double-blind crossover, two hours of diesel exhaust
against filtered air, paired-end RNA-seq on both blood and bronchial brushings
(\textit{Environ Res}). Airway differential expression at FDR $<$ 10\% picks up GPX2 and
NQO1 --- the NRF2 antioxidant arm --- and hallmark mapping gives interferon-$\alpha$,
interferon-$\gamma$, inflammatory and ROS activity in \emph{both} compartments. The
concordance between them was not visible when either compartment was analysed alone. The
authors call it exploratory and it is: one dose, one duration, no exposure--response. But
the question is the right one for panel studies of traffic particles, where bronchoscopy
is the reason sample sizes stay at 32.

\item \textbf{Personal monitoring plus multi-omics narrows PM$_{2.5}$ to four
constituents.} A 313-measurement longitudinal panel in older adults
(\textit{Environ Sci Technol}) screened constituents with a mixed model and three
machine-learning methods and landed on chloride, nickel, barium and selenium. Per IQR,
chloride went with a \textbf{3.48\% (95\% CI 1.50--5.46)} rise in diastolic blood pressure
and nickel with a \textbf{0.09\% (0.05--0.12)} fall in FEV$_1$/FVC, with 212 CpG sites
mediating both cardiovascular and pulmonary endpoints. The constituents are collinear, so
the four winners are correlates and not isolated causes --- but the exposure is
\emph{personal}, which removes the spatial misalignment that usually decides
constituent studies before the statistics start.

\item \textbf{A source that is trivial by mass and dominant by tracer.} Wood-fired pizza
ovens in Milan (\textit{preprint}) contribute 0.085--3.34\,$\mu$g\,m$^{-3}$ of primary
PM$_{10}$ and almost no NO$_x$ --- but \textbf{19.7--34.3\% of levoglucosan} and
\textbf{35.0--37.5\% of benzo(a)pyrene}. Any apportionment that reads levoglucosan as
residential wood burning will mis-assign a third of it in a restaurant-dense city centre.
The ranges are model spreads, not measured intervals, and it is one city.

\item \textbf{The cooking-aerosol record reports its negative result rather than burying
it.} FIGAERO-CIMS on oxidised corn-oil emissions (\textit{Atmos Chem Phys}) gives a
molar-mass-dependent partitioning scheme --- gas-favoured below 120\,g\,mol$^{-1}$,
particle-favoured above 220 --- and then states that a substantial gap remains between
FIGAERO-derived effective volatility and equilibrium thermodynamics, attributing it to
combined measurement and model uncertainty rather than resolving it. Non-equilibrium
partitioning in a source that is a rising share of urban organic aerosol is a direct
problem for any mass-based instrument calibrated against an ambient mix.

\item \textbf{The connector's apparent additions were mostly an artefact, and the DOI gate
caught it.} Of the three connector PMIDs the harvester's fresh set lacked, only one was
genuinely new; the other two had shipped on 8 and 9 September via Crossref carrying
\emph{no PMID}, so a PMID-keyed check missed them. \texttt{check\_dois.py} failed the build
with two \textsc{duplicate-doi} errors. Connector records must be matched on DOI.

\end{itemize}

\clearpage

% ==================================================================== ANALYTICS
\section*{\sffamily\bfseries\color{Deep}Corpus analytics}
\vspace{-1.5mm}{\color{Amber}\rule{\linewidth}{1.1pt}}\vspace{2.5mm}

\noindent\includegraphics[width=\linewidth]{\FIGDIR/f1_subtopics.png}
\figcap{Subtopic assignment of the 10 in-scope records; single-label by dominant
contribution. \textit{Sensing, forecasting \& instrumentation} and \textit{Burden, policy
\& mitigation} tie at 2, and six further clusters carry one apiece --- so 8 of 10 clusters
are represented from only 10 records, and the largest cluster holds just 20\% of the day
--- matching 6 and 7 September for the week's most even spread. \textit{Occupational \& indoor} and \textit{Other clinical endpoints} are the two
empty clusters; the former held 4 records on the backfilled 10 September issue, so the
gap is a one-day artefact rather than a trend.}

\vspace{2mm}
\noindent\includegraphics[width=\linewidth]{\FIGDIR/f2_design_endpoint.png}
\figcap{Left --- study architecture. \textit{Metadata only} is the largest single
architecture at 3, which is an unwelcome thing for it to be; \textit{Chamber /
laboratory}, \textit{Observational --- cohort} and \textit{Review / synthesis} carry 2 each
and \textit{Modelling / inventory} 1. Right --- health endpoint. \textit{Respiratory} leads
at 3 and \textit{No health endpoint} follows at 3 --- the cooking-aerosol, AOD-validation
and pizza-oven records --- with metabolic, cognitive, reproductive and other-clinical
categories carrying one each.}

\vspace{2mm}
\noindent\includegraphics[width=\linewidth]{\FIGDIR/f3_metric_geography.png}
\figcap{Left --- particle metric. \textit{PM (unspeciated) / emissions} leads at 5, which
on this day means three metadata-only records whose particle metric could not be
established plus the two source-characterisation records. PM$_{2.5}$ alone is the exposure
in 2, joint PM$_{2.5}$/PM$_{10}$ in 1, and one record each resolves composition
(the personal-monitoring constituent panel) and optical properties (the AOD product).
Right --- geographic provenance. China contributes 3 and the \textit{Global /
multi-region} fallback 3, of which two are genuinely multi-region reviews and one is a
record whose study region could not be established without an abstract. East Asia
(ex-China) carries 2 --- both Korean --- with Europe and North America one apiece.}

\vspace{2mm}
\noindent\includegraphics[width=\linewidth]{\FIGDIR/f6_lifecourse.png}
\figcap{Life-course windows addressed; counts are non-exclusive and only records with an
identifiable human window are placed. Working-age adulthood and older adulthood carry 2
each --- the Korean asthma cohort and the Vancouver diesel-exhaust volunteers at one end,
the personal-monitoring constituent panel and the 50+ cognition review at the other. The
in-utero window carries 1, the Korean birth-cohort growth study, and it is a
metadata-only record. Childhood and adolescence are both empty.}

\vspace{2mm}
\noindent\includegraphics[width=\linewidth]{\FIGDIR/f4_heatmap.png}
\figcap{Subtopic $\times$ study-architecture cross-tabulation. As on the backfilled
10 September issue, \textbf{every record occupies its own cell} --- 10 records, 10 distinct
subtopic/architecture pairs. Two consecutive days with no cell above one says something
about the window rather than the field: a one- to two-day entry window samples too thinly
for methodological clustering to appear, and any such clustering in this archive is a
property of the rollups, not of a daily issue. The sensing row again carries no human
health endpoint in either of its two records.}

\vspace{2mm}
\noindent\includegraphics[width=\linewidth]{\FIGDIR/f5_forest.png}
\figcap{All four quantitative estimates in the corpus on a common logarithmic axis.
\textbf{All four come from one study}, the Korean nationwide asthma cohort, and the axis
should be read as a within-study comparison rather than a synthesis. The top two are the
main exposure contrasts, per 10\,$\mu$g\,m$^{-3}$ of PM$_{2.5}$ and of PM$_{10}$; the
bottom two are the extremes of the three-way stratification on age, Medical Aid status and
urbanicity, and the gap between 1.16 and 1.80 is the paper's actual claim. Intervals are
narrow enough that they are barely visible at this scale --- that is a sample-size artefact
of 8.25 million people and not a precision claim, and the Signal of the day argues that
the dominant uncertainty here is not sampling error. Estimates omitted: the
personal-monitoring constituent results are percentage changes in blood pressure and
FEV$_1$/FVC rather than ratios and cannot be placed on this axis.}

\clearpage

% ==================================================================== DIGEST
\section*{\sffamily\bfseries\color{Deep}Paper-level digest}
\vspace{-1.5mm}{\color{Amber}\rule{\linewidth}{1.1pt}}\vspace{1mm}

{\fontsize{7.4}{9}\selectfont\color{Slate}Ordered by cluster. Each entry gives the
methodological core and the finding that matters, not an abstract paraphrase. Entries
marked \textsc{metadata only} had no abstract in existence on the entry date; for these
the entry states what the record establishes and what it does not, and asserts no
finding.\par}

\band{Neuro / mental health\hfill 1 record}{Deep}

\paperentry{Older-adult cognition review 2026}
{Air pollution and cognitive outcomes in older adults: a systematic review}
{Rev Environ Health}
{PubMed searched to February 2025 for observational studies from 2000 onward on ambient air pollution and cognition in adults 50+; 36 studies included. PM$_{2.5}$ is the most-studied pollutant, and most studies report worse global or domain-specific cognition, faster decline, or higher impairment and dementia risk with higher PM$_{2.5}$. Other pollutants trend adverse; ozone is the least consistent. Heterogeneity is attributed to exposure-assessment method, cognitive instrument, design and confounding control. Two limitations bound what this can support: a single-database search with no quantitative pooling and no formal risk-of-bias tool, so the conclusion is a vote count, not a weighted synthesis, and publication bias in a literature this size is unaddressed. The exposure-assessment heterogeneity it names is the actionable part}
{10.1515/reveh-2026-0022}

\band{Cardiovascular \& metabolic\hfill 1 record}{Teal}

\paperentry{Constituent multiomics cardiopulmonary 2026}
{PM$_{2.5}$ Constituents and Cardiopulmonary Dysfunctions in Older Adults: A Multiomics Approach Unveiling Molecular Pathways}
{Environ Sci Technol}
{Longitudinal panel, 313 repeated measurements in older adults, personal rather than ambient PM$_{2.5}$ with constituent speciation, screened by a linear mixed model plus three machine-learning methods. Chloride, nickel, barium and selenium emerged as the key constituents. Per IQR, chloride was associated with a 3.48\% (95\% CI 1.50-5.46) rise in diastolic blood pressure and nickel with a 0.09\% (0.05-0.12) fall in FEV1/FVC. Mediation implicated 212 CpG sites acting on both cardiovascular and pulmonary endpoints and 7 genes spanning more than two constituent-function pairs, with 22 CpGs linked by eQTM to 14 coding genes including CERS6 and LRCH3. Limitations: 313 measurements is a small panel for a screen over this many constituents and CpGs, the constituent set is collinear so the four winners are correlates rather than isolated causes, and the mediation is cross-sectional within visit. Personal monitoring is the strength - it removes the spatial misalignment that defeats constituent studies built on central-site speciation}
{10.1021/acs.est.6c04206}

\band{Respiratory \& allergic\hfill 1 record}{Sky}

\paperentry{Korea 8.25M asthma cohort 2026}
{Cumulative Vulnerability to Particulate Matter and Incident Asthma: Joint Effects of Urbanicity, Poverty, and Age in a Nationwide Cohort}
{Allergy}
{8,252,498 adults in the Korean National Health Insurance Service database linked to 1-km long-term PM estimates, Cox models for incident asthma. Per 10 $\mu$g\,m$^{-3}$, PM$_{2.5}$ gave HR 1.33 (95\% CI 1.32-1.34) and PM$_{10}$ HR 1.38 (1.37-1.38). Stratification by age, Medical Aid status and urbanicity produced a monotone gradient from HR 1.16 (1.14-1.18) in younger rural adults without Medical Aid to 1.80 (1.66-1.94) in older urban Medical Aid recipients, p for interaction < 0.001. Limitation: incident asthma in claims data is diagnosis-coded, and urban Medical Aid recipients have higher healthcare contact, so part of the 1.16-to-1.80 gradient is plausibly ascertainment rather than susceptibility - the confidence intervals are narrow enough to make that distinction matter. Exposure contrast at 1 km is coarse relative to the urbanicity modifier being tested, which is the resolution argument for dense low-cost networks}
{10.1111/all.70502}

\band{Reproductive \& developmental\hfill 1 record}{Sage}

\paperentry{Korean birth-cohort growth mixture 2026}
{Prenatal bisphenols and ambient air pollutants are differentially associated with early childhood length and weight from birth to 36 months: a multi-domain mixture analysis in a Korean birth cohort}
{Environ Pollut}
{\textsc{metadata only.} No abstract was available on the entry date from Crossref per-DOI, Europe PMC, PubMed efetch or Semantic Scholar. Title establishes a Korean birth cohort with a multi-domain mixture analysis placing prenatal bisphenols alongside ambient air pollutants against length and weight trajectories to 36 months, and states that the two exposure domains behave differently - but not in which direction, for which pollutant, or at what magnitude. Nothing is asserted here and the record is tier C. Flag for re-reading: multi-domain mixture methods that carry chemical and air exposures in one model are the methodological question this archive should track}
{10.1016/j.envpol.2026.129145}

\band{Mechanistic toxicology\hfill 1 record}{Violet}

\paperentry{Vancouver DE blood-lung concordance 2026}
{Concordant transcriptomics analysis of the lungs and blood after exposure to diesel exhaust: a randomized, controlled human exposure study}
{Environ Res}
{Double-blind crossover controlled human exposure, 32 participants, two hours of diesel exhaust and of filtered air on separate occasions, with paired-end RNA-seq on blood and bronchial brushings. Airway DEGs at FDR < 10\% include upregulated GPX2 and NQO1, i.e. the NRF2 antioxidant arm. Hallmark pathway mapping gives IFN-alpha, IFN-gamma, inflammatory response and ROS activity in both compartments, and a concordant analysis finds correlated blood-airway expression that neither compartment showed alone. The authors are explicit that this is exploratory. n = 32 at a single two-hour dose, no exposure-response, and the concordance is correlational - but it is the right design question: whether a blood transcriptomic readout can stand in for bronchoscopy in panel studies of traffic-related particles}
{10.1016/j.envres.2026.125625}

\band{Exposure assessment \& modelling\hfill 1 record}{Clay}

\paperentry{Milan pizza-oven dispersion 2026}
{Assessment of Air Quality Impacts from Wood-Fired Pizza Ovens: Comparison with Major Urban Emission Sources}
{Preprints.org}
{PREPRINT. Milan. Measured wood-fired-oven emission factors feed a dispersion model, with road traffic and domestic heating simulated on the same framework and all three compared against ARPA Lombardia monitors. Oven contributions concentrate in the city centre: 0.085-3.34 $\mu$g\,m$^{-3}$ of primary PM$_{10}$, only 0.0044-1.75 $\mu$g\,m$^{-3}$ of NOx, but 19.70-34.28\% of levoglucosan and 34.95-37.47\% of benzo(a)pyrene. That split is the finding - a source that is minor by mass and major by tracer, which is exactly the pattern that breaks levoglucosan-based residential-wood-burning apportionment in a restaurant-dense centre. Limitations: preprint, single city, emission factors from a small appliance sample, and the ranges are model spreads rather than measured confidence intervals}
{10.20944/preprints202609.0792.v1}

\band{Sensing, forecasting \& instrumentation\hfill 2 records}{Amber}

\paperentry{Cooking aerosol FIGAERO partitioning 2026}
{Unravelling gas-particle partitioning dynamics in cooking aerosol oxidation through FIGAERO-CIMS analysis}
{Atmos Chem Phys}
{Heated corn oil emissions aged in a Go:PAM oxidation flow reactor, with FIGAERO-I-CIMS and a 2-D VBS parameterisation tracking composition, volatility and partitioning. Kinetic calculations give molar-mass-dependent partitioning with the surface process rate-limiting. Primary cooking aerosol sits biased toward the gas phase relative to equilibrium; in secondary aerosol the crossover is near 120-220 g\,mol$^{-1}$, below which compounds favour gas and above which they favour particle. The honest result is the negative one: a substantial unresolved gap remains between FIGAERO-derived effective volatility and equilibrium thermodynamic estimates, which the authors attribute to combined measurement and model uncertainty rather than resolving. Cooking is a rising share of urban OA as combustion sources are controlled, and non-equilibrium partitioning is a direct problem for any mass-based sensor calibrated on ambient mix}
{10.5194/acp-26-12793-2026}

\paperentry{FY-4B/AGRI AOD evaluation 2026}
{Evaluation of the FY-4B/AGRI Land Aerosol Optical Depth Product and Comparison with Geostationary Satellite Products over Asia}
{Atmos Environ}
{\textsc{metadata only.} No abstract was available on the entry date from any of the four sources queried - the same same-day Elsevier deposit lag logged repeatedly since 2026-09-08. Title and journal establish a validation of the FY-4B/AGRI land AOD product against other geostationary retrievals over Asia. No validation statistics, no reference network and no spatial or temporal coverage are available, so nothing is asserted here and the record is tier C. Worth re-reading when the abstract deposits: geostationary AOD is the only route to sub-daily satellite-surface PM fusion, and product-versus-product agreement is the term that dominates the error budget of any such fusion}
{10.1016/j.atmosenv.2026.122354}

\band{Burden, policy \& mitigation\hfill 2 records}{Coral}

\paperentry{China PM$_{1}$ burden projection 2026}
{Ambient PM$_{1}$-Related Health Burdens and Mortality Influencing Factors in China: Historical Trends to Future Projections}
{Environ Pollut}
{\textsc{metadata only.} No abstract was available on the entry date from any of the four sources queried. Title establishes a China-wide PM$_{1}$ attributable-burden analysis running historical trends into projections with a decomposition of mortality drivers. No exposure product, no concentration-response function, no scenario definition and no burden figure are available, so nothing is asserted and the record is tier C. Worth re-reading: PM$_{1}$ burden work is scarce because almost no regulatory network measures PM$_{1}$, which makes whatever exposure surface this study used the interesting part}
{10.1016/j.envpol.2026.129144}

\paperentry{COVID air-quality lessons review 2026}
{COVID-19 Pandemic as an Environmental Health Indicator: Lessons Learned from Air Pollution Changes and Climate Implications}
{Preprints.org}
{PREPRINT, single author. Reviews the lockdown natural experiment: short-term outdoor air-quality improvement and greenhouse-gas reduction across many regions, the reported interactions between meteorology, pollutant concentrations and respiratory viral transmission, and the argument that temporary activity reductions demonstrate the ceiling of behavioural mitigation rather than a strategy. No systematic search, no inclusion criteria, no quantitative synthesis, not peer reviewed, and the pollution-COVID-transmission literature it draws on is itself heavily confounded by the timing of interventions. Logged for completeness of the burden-and-policy subtopic; carries no estimate this archive should propagate}
{10.20944/preprints202609.0712.v1}


\clearpage

% ==================================================================== TABLE
\section*{\sffamily\bfseries\color{Deep}Corpus register}
\vspace{-1.5mm}{\color{Amber}\rule{\linewidth}{1.1pt}}\vspace{2.5mm}

{\fontsize{6.7}{8.4}\selectfont
\setlength{\tabcolsep}{2.4pt}
\begin{longtable}{@{}>{\RaggedRight}p{25mm}>{\RaggedRight}p{41mm}>{\RaggedRight}p{27mm}%
>{\RaggedRight}p{23mm}>{\RaggedRight}p{18mm}>{\RaggedRight}p{15mm}@{}}
\toprule
\textsf{\textbf{First author}} & \textsf{\textbf{Focus}} & \textsf{\textbf{Journal}} &
\textsf{\textbf{Design}} & \textsf{\textbf{Metric}} & \textsf{\textbf{Region}}\\
\midrule
\endfirsthead
\toprule
\textsf{\textbf{First author}} & \textsf{\textbf{Focus}} & \textsf{\textbf{Journal}} &
\textsf{\textbf{Design}} & \textsf{\textbf{Metric}} & \textsf{\textbf{Region}}\\
\midrule
\endhead
\multicolumn{6}{@{}l@{}}{\cellcolor{Mist}\textsf{\textbf{\textcolor{Deep}{Neuro / mental health}}}}\\[0.6mm]
\href{https://doi.org/10.1001/jamanetworkopen.2026.25015}{Dales} & Wildfire-sourced PM$_{2.5}$ and ED visits for migraine and other primary headache & \textit{JAMA Netw Open} & Case-crossover & PM$_{2.5}$ (wildfire) & Canada\\
\href{https://doi.org/10.1093/ageing/afag215}{Glover} & Umbrella meta-review of urban environmental exposures and cognitive health & \textit{Age Ageing} & Umbrella review / meta-analysis & PM$_{2.5}$ & Global\\
\href{https://doi.org/10.1016/j.nsa.2026.107019}{Goossens} & In utero air pollution exposure and pre-/postnatal brain development & \textit{Neurosci Appl} & Narrative review & PM$_{2.5}$ / PM$_{10}$ / PM$_{\mathrm{coarse}}$ & Global\\
\href{https://doi.org/10.1016/j.isci.2026.116617}{Ma} & PM$_{2.5}$-induced neurotoxicity and depression: exposure risks and mechanisms & \textit{iScience} & Narrative review & PM$_{2.5}$ + components & Global\\
\href{https://doi.org/10.3760/cma.j.cn112150-20250804-00752}{Pan} & Maternal PM exposure and offspring neurodevelopment: effects and mechanisms & \textit{Zhonghua Yu Fang Yi Xue Za Zhi} & Narrative review & PM (bulk) & China\\
\addlinespace[0.9mm]
\multicolumn{6}{@{}l@{}}{\cellcolor{Mist}\textsf{\textbf{\textcolor{Teal}{Cardiovascular \& metabolic}}}}\\[0.6mm]
\href{https://doi.org/10.3760/cma.j.cn112338-20251120-00843}{Huang} & PM$_{2.5}$ and dyslipidaemia in rural Chinese children; dietary effect modification & \textit{Zhonghua Liu Xing Bing Xue Za Zhi} & Cross-sectional (multistage) & PM$_{2.5}$ & China\\
\href{https://doi.org/10.32394/pe/218966}{Kazmierska} & Air pollution and cardiovascular disease: review of current data & \textit{Przegl Epidemiol} & Narrative review & PM$_{2.5}$ / PM$_{10}$ & Global\\
\href{https://doi.org/10.1007/s00125-026-06814-2}{Xi} & Air pollution, built environment and temperature vs. prediabetes phenotypes & \textit{Diabetologia} & Prospective cohort & PM$_{2.5}$ & Germany\\
\href{https://doi.org/10.1161/JAHA.125.046917}{Yang} & Short-term PM$_{2.5}$ and hospitalisation/ED use in atrial fibrillation & \textit{J Am Heart Assoc} & Case-crossover & PM$_{2.5}$ (wildfire-influenced) & USA\\
\addlinespace[0.9mm]
\multicolumn{6}{@{}l@{}}{\cellcolor{Mist}\textsf{\textbf{\textcolor{Sky}{Reproductive \& developmental}}}}\\[0.6mm]
\href{https://doi.org/10.1016/j.jes.2025.10.043}{Chen} & Prenatal air pollution, fetal growth trajectory phenotypes and child neurodevelopment & \textit{J Environ Sci (China)} & Prospective cohort & PM$_{2.5}$ / PM$_{10}$ & China\\
\href{https://doi.org/10.1016/j.ecoenv.2026.120566}{Huang} & Joint prenatal bisphenol + air-pollutant mixture exposure and preschool IQ & \textit{Ecotoxicol Environ Saf} & Prospective cohort (mixtures) & PM$_{2.5}$ / PM$_{10}$ & China\\
\href{https://doi.org/10.1007/s40572-026-00552-8}{Peterson} & Ambient air pollution, PAHs and pubertal development: critical review & \textit{Curr Environ Health Rep} & Critical review & PM$_{2.5}$ components & Global\\
\href{https://doi.org/10.1016/j.envres.2026.125313}{Xie} & Long-term PM$_{2.5}$ and timing of pubertal development in boys and girls & \textit{Environ Res} & Prospective cohort & PM$_{2.5}$ & China\\
\addlinespace[0.9mm]
\multicolumn{6}{@{}l@{}}{\cellcolor{Mist}\textsf{\textbf{\textcolor{Sage}{Respiratory \& allergic}}}}\\[0.6mm]
\href{https://doi.org/10.1007/s00484-026-03289-z}{Bayar Muluk} & Climate change and otolaryngology: airway disease impacts and emission reduction & \textit{Int J Biometeorol} & Narrative review & PM$_{2.5}$ & Global\\
\href{https://doi.org/10.4103/mgr.MEDGASRES-D-25-00164}{Liao} & Short-term air pollution, respiratory symptoms and lung function in asthma & \textit{Med Gas Res} & Panel study & PM$_{2.5}$ / PM$_{10}$ & China\\
\href{https://doi.org/10.3390/epidemiologia7040097}{Sanchez de Prada} & Ambient PM and asthma-coded admissions in children 0-3 y, Spain & \textit{Epidemiologia} & Ecological panel & PM$_{2.5}$ / PM$_{10}$ & Spain\\
\href{https://doi.org/10.1016/j.jhazmat.2026.143076}{Yang} & Allergenic vs. toxic risk of airborne seed hairs compared with ambient PM$_{2.5}$ & \textit{J Hazard Mater} & Field measurement & PM$_{2.5}$ + bioaerosol & China\\
\addlinespace[0.9mm]
\multicolumn{6}{@{}l@{}}{\cellcolor{Mist}\textsf{\textbf{\textcolor{Violet}{Mechanistic toxicology}}}}\\[0.6mm]
\href{https://doi.org/10.1016/j.jes.2025.12.074}{Ding} & PM$_{2.5}$ promotes vascular calcification via miR-27b-5p/HES1 signalling & \textit{J Environ Sci (China)} & In vitro + organoid + murine & PM$_{2.5}$ & China\\
\href{https://doi.org/10.1007/s00360-026-01698-5}{Maraschi} & Osmoregulatory disruption and metal bioaccumulation from settleable atmospheric PM & \textit{J Comp Physiol B} & Controlled exposure (ecotox) & Settleable APM & Brazil\\
\href{https://doi.org/10.1016/j.jes.2025.11.018}{Wu} & Macrophage NLRP3 inflammasome in real-ambient PM-induced hepatic glucose dysregulation & \textit{J Environ Sci (China)} & Real-ambient murine + co-culture & Real-ambient PM & China\\
\addlinespace[0.9mm]
\multicolumn{6}{@{}l@{}}{\cellcolor{Mist}\textsf{\textbf{\textcolor{Clay}{Occupational \& indoor}}}}\\[0.6mm]
\href{https://doi.org/10.3389/fmedt.2026.1796828}{Chen} & Nested suction instrument to cut intraoperative surgical-smoke PM$_{2.5}$ & \textit{Front Med Technol} & Bench / device experiment & PM$_{2.5}$ (surgical smoke) & China\\
\href{https://doi.org/10.1093/annweh/wxag055}{Cole} & Determinants of respirable crystalline silica exposure in Australian tunnelling & \textit{Ann Work Expo Health} & Observational exposure study & Respirable dust / RCS & Australia\\
\href{https://doi.org/10.1016/j.jhazmat.2026.143081}{Ding} & Stage-resolved hair Zn-isotope dynamics as a firefighting combustion-exposure tracer & \textit{J Hazard Mater} & Repeated-measures biomarker & Combustion-derived PM & China\\
\href{https://doi.org/10.3390/jox16040119}{Ramsden} & Xenobiotic hazards in aircraft cabin air & \textit{J Xenobiot} & Review + risk economics & Ultrafine particles & Global\\
\addlinespace[0.9mm]
\multicolumn{6}{@{}l@{}}{\cellcolor{Mist}\textsf{\textbf{\textcolor{Amber}{Sensing, forecasting \& instrumentation}}}}\\[0.6mm]
\href{https://doi.org/10.1016/j.jenvman.2026.130515}{Shi} & Adaptive spatiotemporal graph transformer for multi-site multi-step PM$_{2.5}$ forecasting & \textit{J Environ Manage} & Deep-learning model & PM$_{2.5}$ & China / India\\
\href{https://doi.org/10.3390/bios16070380}{Tsuchiya} & Seasonal and indoor-outdoor characteristics of fluorescent aerosol particles in offices & \textit{Biosensors} & Real-time sensor campaign & Bioaerosol / coarse & Japan\\
\addlinespace[0.9mm]
\multicolumn{6}{@{}l@{}}{\cellcolor{Mist}\textsf{\textbf{\textcolor{Amber}{Exposure assessment \& modelling}}}}\\[0.6mm]
\href{https://doi.org/10.1029/2025GH001621}{Solorzano-Araque} & Multiscale exposure modelling: LOTOS-EUROS + mobility surveys + particle toxicity & \textit{GeoHealth} & CTM + mobility integration & PM$_{2.5}$ / PM$_{10}$ & Colombia\\
\addlinespace[0.9mm]
\multicolumn{6}{@{}l@{}}{\cellcolor{Mist}\textsf{\textbf{\textcolor{Coral}{Burden, policy \& mitigation}}}}\\[0.6mm]
\href{https://doi.org/10.1002/hsr2.72886}{Karamzad} & Ischaemic heart disease burden in MENA 1990-2021; PM attributable fraction 28.2\% & \textit{Health Sci Rep} & GBD secondary analysis & Ambient PM & MENA\\
\href{https://doi.org/10.1371/journal.pone.0347989}{Rees} & HEPA filtration to reduce respiratory infection in care homes & \textit{PLoS One} & RCT process evaluation & Filtered PM & UK\\
\href{https://doi.org/10.1016/j.jes.2026.01.009}{Shan} & Policy-driven emission reductions and post-pandemic rebound, high-resolution inventories & \textit{J Environ Sci (China)} & Emission inventory & PM emissions & China\\
\addlinespace[0.9mm]
\multicolumn{6}{@{}l@{}}{\cellcolor{Mist}\textsf{\textbf{\textcolor{Slate}{Other clinical endpoints}}}}\\[0.6mm]
\href{https://doi.org/10.3390/dj14070440}{Lupita} & PM$_{2.5}$, preventive-care capacity and caries/edentulism burden across the EU & \textit{Dent J} & Ecological panel & PM$_{2.5}$ & EU-27\\
\href{https://doi.org/10.1016/j.jdin.2026.06.008}{Park} & Long-term PM exposure and risk of androgenetic alopecia & \textit{JAAD Int} & Nationwide cohort & PM$_{2.5}$ / PM$_{10}$ & South Korea\\
\href{https://doi.org/10.3171/2026.3.PEDS2615}{Sun} & Environmental exposure and area deprivation vs. paediatric hydrocephalus outcomes & \textit{J Neurosurg Pediatr} & Registry cohort + geospatial & PM$_{2.5}$ (satellite-derived) & USA\\
\addlinespace[0.9mm]
\bottomrule
\end{longtable}}


\vspace{4mm}
\begin{keybox}
{\sffamily\bfseries\footnotesize\color{Deep}Provenance and method}\\[1.2mm]
{\fontsize{7.2}{8.8}\selectfont
\textbf{Window.} Entry dates 10--11 September 2026. The harvester computed the window from
\texttt{last\_entry\_date} before the backfilled 10 September issue had been saved, so the
window deliberately overlaps that day; the overlap is absorbed by \texttt{seen.json} and a
direct re-index of the corpus, which together dropped 28 of the 95 unique records as
already archived. No record appears in both issues.\par
\vspace{0.8mm}
\textbf{Sources.} \textbf{PubMed} E-utilities on \texttt{[Date - Entry]}, health and
instrumentation axes as separate queries (12 hits and 1); the \textbf{PubMed connector} run
independently returned 14 on the health axis. \textbf{Europe PMC}
\texttt{CREATION\_DATE} (27 hits). \textbf{Crossref by ISSN} across the watch's journal
list (64 hits) --- the only leg reaching \textit{Atmos.\ Chem.\ Phys.}, which supplied the
cooking-aerosol record. \textbf{arXiv} relevance query with a local recency screen ---
nothing passed. \textbf{OpenAlex} HTTP 429 again (rate limit, not the 403 the source note
still records); that leg has been degraded since July. \textbf{Consensus} returned six hits
across a low-cost-sensor calibration/drift sweep and a satellite--surface fusion sweep;
all six are 2019--2025 archive papers with Crossref \emph{created} dates far outside this
window, so the leg is a topical nil for this issue rather than a contribution.\par
\vspace{0.8mm}
\textbf{Connector reconciliation.} Of the connector's 14 health-axis PMIDs, three were
absent from the harvester's fresh set. One is genuinely new and ships (the Vancouver
diesel-exhaust crossover). The other two had already shipped on 8 and 9 September, having
arrived then via Crossref with an empty PMID field, so a PMID-keyed screening check could
not see them. \texttt{check\_dois.py} rejected the build with two \textsc{duplicate-doi}
failures and both records were removed. The screening check for connector-surfaced records
is being moved onto DOI.\par
\vspace{0.8mm}
\textbf{Screening.} 105 raw records, 95 after cross-source deduplication on
PMID\,$\rightarrow$\,DOI\,$\rightarrow$\,normalised title, 28 already archived, 67 fresh.
\textbf{10 in scope, 28 newly excluded} --- 27 from the harvest and one connector-only
record --- each with a logged reason in \texttt{state/rejected.jsonl}. A further 32
exclusions were suppressed as duplicates of rows already written under 10 September, so
each exclusion is recorded once in the ledger.\par
\vspace{0.8mm}
\textbf{Metadata-only records.} Three records had no abstract in existence anywhere on the
entry date --- checked per-DOI against Crossref, Europe PMC, PubMed \texttt{efetch} and
Semantic Scholar. All three are unambiguously in scope on the title and ship as
\textsc{metadata only}, tier C, asserting no finding. This has recurred on every issue
since 8 September and is treated as expected behaviour for same-day publisher deposits.\par
\vspace{0.8mm}
\textbf{Trial watch.} ClinicalTrials.gov returned \textbf{0} PM or air-pollution
registrations or updates across 10--11 September against \textbf{2{,}269 registry-wide} on
the no-term control --- a topical zero, not an outage. Accumulated trial watch surfaces in
the weekly rollup (next: Saturday 12 September).\par
\vspace{0.8mm}
\textbf{Labels and estimates.} Subtopic, design, particle-metric and geography labels are
assigned from title, abstract, keywords and MeSH; assignment is single-label by dominant
contribution and is a judgement, not a controlled vocabulary. Effect estimates are
transcribed verbatim from abstracts --- see the forest-plot caption, which records which
estimates are omitted and why. All DOIs resolve to the publisher of record and are checked
automatically before build. Abstracts remain the copyright of their publishers.\par}
\end{keybox}

\end{document}
